\section{Carregamento, linking e inicialização}
\label{sec:class-loading}

O cap.~5 da JVMS divide a vida de uma classe em três fases.

\begin{description}
\item[Carregamento (§5.3)] Um \termo{class loader}{Objeto que encontra os
  bytes de uma classe e a define na JVM.} encontra os bytes da classe.
  Há dois tipos: o \emph{bootstrap class loader}, da própria JVM, e os
  definidos pelo usuário, subclasses de \texttt{ClassLoader}. Um loader
  pode definir a classe ou delegar a outro.
\item[Linking (§5.4)] \emph{Verificação} confere que o bytecode é
  estruturalmente correto (falha: \texttt{VerifyError}).
  \emph{Preparação} cria os campos \texttt{static} com valores padrão,
  sem rodar código Java. \emph{Resolução} transforma referências
  simbólicas do \termo{constant pool}{Tabela de constantes do
  \texttt{.class}: nomes de classes, métodos, campos, strings e números.
  Ver a seção~\ref{sec:constant-pool}.} em referências concretas; a
  especificação permite fazer isso de forma preguiçosa, só quando a
  referência é usada.
\item[Inicialização (§5.5)] Roda o inicializador da classe
  (\texttt{<clinit>}: atribuições a campos \texttt{static} e blocos
  \texttt{static}). Acontece, por exemplo, na primeira execução de
  \texttt{new}, \texttt{getstatic}, \texttt{putstatic} ou
  \texttt{invokestatic} que referencia a classe.
\end{description}

\paragraph{O bug que isso explica.}
Se o bloco \texttt{static} de uma classe lança uma exceção (uma
configuração que não existe, um \texttt{Integer.parseInt} de variável de
ambiente vazia), o primeiro uso recebe \texttt{ExceptionInInitializerError}
com a causa original. A JVMS (§5.5, passos 11 e 12) manda marcar a classe
como \emph{erroneous}; a partir daí, todo uso recebe
\texttt{NoClassDefFoundError} (passo 5), que parece um problema de
classpath e não é. Quando vir \texttt{NoClassDefFoundError} num serviço
que funcionava, procure o primeiro \texttt{ExceptionInInitializerError}
no log.

\paragraph{A decisão que isso muda.}
Inicialização \texttt{static} pesada (abrir conexão, ler arquivo) roda no
primeiro uso, na thread que chegou primeiro, e transforma qualquer falha
nesse par de erros. Prefira inicialização explícita, num componente com
ciclo de vida (um bean), em que o erro aparece na subida.

\paragraph{Startup e o cache AOT.}
Carregar e ligar milhares de classes custa tempo a cada subida. O
\citetitle{doc:jep483} (JDK~24) guarda classes já carregadas e ligadas
num cache \termo{AOT}{\emph{Ahead-of-time}: compilação ou preparação
feita antes da execução, em oposição ao JIT, que compila durante a
execução.} gerado por uma execução de treino; o \citetitle{doc:jep514}
(JDK~25) reduz a criação do cache a um comando
(\texttt{java -XX:AOTCacheOutput=app.aot \ldots}), e o
\citetitle{doc:jep515} (JDK~25) guarda também perfis de métodos, para o
JIT começar mais cedo. O JEP~483 informa, sem dizer quem mediu nem em que
máquina, que o Spring PetClinic 3.2.0 subiu em 4,486~s no JDK~23 e em
2,604~s no JDK~24 com o cache (42\% mais rápido). Trate como número do
próprio projeto OpenJDK, não como garantia para o seu serviço.

\paragraph{Ponte com Quarkus.}
O guia \citetitle{doc:quarkus-class-loading} descreve uma hierarquia
própria em dev mode: um \emph{Base Runtime ClassLoader} persistente para
as dependências e um \emph{Runtime ClassLoader} para as classes da
aplicação, recriado quando a aplicação reinicia; é isso que permite o
live reload sem reiniciar a JVM. No pacote padrão de produção
(\texttt{fast-jar}), um loader que indexa as classes em tempo de build;
em executável nativo, classpath plano. A Parte~\ref{parte:quarkus} volta a isso.
